% Preamble of the WFDB-x user's manual (wfdb-x-user-manual.tex).
% Built with pdflatex (latexmk; see README.md). The source is ASCII: special characters are
% written as LaTeX commands.

% ---- Engine, encoding and fonts ----------------------------------------------------------
\usepackage{iftex}
\ifPDFTeX
  \usepackage[T1]{fontenc}
  \usepackage[utf8]{inputenc}
\else
  \usepackage{fontspec}
\fi
\usepackage{XCharter}                               % text: Charter
\IfFileExists{sourcesans.sty}                       % headings, labels: Source Sans (\sffamily)
  {\usepackage{sourcesans}}{\usepackage{sourcesanspro}}  % (older TeX distributions)
\usepackage[varqu,varl,scaled=0.9]{zi4}             % code: Inconsolata, upright quotes, bold
\usepackage{textcomp}
\usepackage{amsmath,amssymb}
\usepackage{microtype}
\DisableLigatures[-,<,>]{family=tt*}           % "--", "<<", ">>" stay as typed in code

% ---- Colours (colour-blind-safe; red is never used: it is for alarms) --------------------
\usepackage[table]{xcolor}
\definecolor{titlecolor}{HTML}{12324A}               % headings, title page
\definecolor{accent}{HTML}{0072B2}                   % Okabe-Ito blue: numbers, links, rules
\definecolor{okskyblue}{HTML}{56B4E9}
\definecolor{okgreen}{HTML}{009E73}
\definecolor{okorange}{HTML}{E69F00}
\definecolor{okviolet}{HTML}{7B3294}                 % identifying data (PHI)
\definecolor{rulecolor}{HTML}{B8C4CE}
\definecolor{mutedtext}{HTML}{5A6B7B}
\definecolor{codebg}{HTML}{F5F7F9}
\definecolor{codecomment}{HTML}{5A6B7B}
\definecolor{codekeyword}{HTML}{005A8C}
\definecolor{codestring}{HTML}{00785A}
\definecolor{noteback}{HTML}{EBF4FA}
\definecolor{phiback}{HTML}{F4EDF8}
\definecolor{tipback}{HTML}{E9F6F1}
\definecolor{tableheadback}{HTML}{EEF2F5}

% ---- Page --------------------------------------------------------------------------------
\usepackage[a4paper,top=28mm,bottom=27mm,inner=27mm,outer=25mm,headheight=15pt,
  headsep=9mm,footskip=11mm]{geometry}
\usepackage[skip=0.55em plus 0.2em minus 0.1em]{parskip}
\raggedbottom
\clubpenalty=10000
\widowpenalty=10000
\setcounter{tocdepth}{2}
\setcounter{secnumdepth}{3}

% ---- Headings and the table of contents --------------------------------------------------
\usepackage{titlesec}
\usepackage{titletoc}
\titleformat{\part}[display]
  {\normalfont\sffamily\centering\color{titlecolor}}
  {\large\bfseries\color{accent}\MakeUppercase{\partname}\enspace\thepart}
  {14pt}
  {\Huge\bfseries}
\titleformat{\chapter}[display]
  {\normalfont\sffamily\color{titlecolor}}
  {\large\bfseries\color{accent}\MakeUppercase{\chaptertitlename}\enspace\thechapter}
  {4pt}
  {{\color{accent}\titlerule[1.2pt]}\vspace{10pt}\Huge\bfseries}
\titleformat{name=\chapter,numberless}[display]
  {\normalfont\sffamily\color{titlecolor}}
  {}
  {0pt}
  {{\color{accent}\titlerule[1.2pt]}\vspace{10pt}\Huge\bfseries}
\titlespacing*{\chapter}{0pt}{0pt}{26pt}
\titleformat{\section}
  {\normalfont\sffamily\Large\bfseries\color{titlecolor}}{\color{accent}\thesection}{0.75em}{}
\titleformat{\subsection}
  {\normalfont\sffamily\large\bfseries\color{titlecolor}}{\color{accent}\thesubsection}{0.75em}{}
\titleformat{\subsubsection}
  {\normalfont\sffamily\normalsize\bfseries\color{titlecolor}}{\thesubsubsection}{0.75em}{}
\titleformat{\paragraph}[runin]
  {\normalfont\sffamily\normalsize\bfseries\color{titlecolor}}{}{0pt}{}
\titlespacing*{\section}{0pt}{1.6em plus 0.4em minus 0.2em}{0.5em}
\titlespacing*{\subsection}{0pt}{1.3em plus 0.3em minus 0.2em}{0.4em}
\titlespacing*{\subsubsection}{0pt}{1.1em plus 0.3em minus 0.2em}{0.3em}
\titlespacing*{\paragraph}{0pt}{0.8em plus 0.2em}{0.6em}

\titlecontents{part}[0pt]
  {\addvspace{1.3em}\sffamily\bfseries\large\color{titlecolor}}
  {\partname\ \thecontentslabel\quad}{}{\hfill\contentspage}[\addvspace{0.3em}]
\titlecontents{chapter}[2.2em]
  {\addvspace{0.55em}\sffamily\bfseries}
  {\contentslabel{2.2em}}{\hspace*{-2.2em}}{\hfill\contentspage}
\titlecontents{section}[4.9em]
  {\small}{\contentslabel{2.7em}}{\hspace*{-2.7em}}
  {\titlerule*[0.7pc]{.}\contentspage}
\titlecontents{subsection}[8.4em]
  {\small\color{mutedtext}}{\contentslabel{3.5em}}{\hspace*{-3.5em}}
  {\titlerule*[0.7pc]{.}\contentspage}
\titlecontents{figure}[2.6em]
  {\small}{\contentslabel{2.6em}}{\hspace*{-2.6em}}{\titlerule*[0.7pc]{.}\contentspage}

% ---- Running heads -----------------------------------------------------------------------
\usepackage{fancyhdr}
\pagestyle{fancy}
\fancyhf{}
\fancyhead[L]{\sffamily\small\color{mutedtext}\nouppercase{\leftmark}}
\fancyhead[R]{\sffamily\small\bfseries\color{titlecolor}\thepage}
\fancyfoot[L]{\sffamily\footnotesize\color{mutedtext}WFDB-x\enspace\textbar\enspace
  User's manual}
\fancyfoot[R]{\sffamily\footnotesize\color{mutedtext}Version \wfdbxversion}
\renewcommand{\headrule}{{\color{rulecolor}\hrule height 0.4pt width\headwidth}}
\renewcommand{\footrulewidth}{0pt}
\renewcommand{\chaptermark}[1]{\markboth{%
  \ifnum\value{secnumdepth}>-1 \ifnum\value{chapter}>0 \thechapter\quad\fi\fi #1}{}}
\fancypagestyle{plain}{\fancyhf{}%
  \fancyfoot[L]{\sffamily\footnotesize\color{mutedtext}WFDB-x\enspace\textbar\enspace
    User's manual}%
  \fancyfoot[R]{\sffamily\footnotesize\bfseries\color{titlecolor}\thepage}%
  \renewcommand{\headrulewidth}{0pt}}

% ---- Lists, figures, tables --------------------------------------------------------------
\usepackage{enumitem}
\setlist{itemsep=0.15em,topsep=0.3em,parsep=0.1em}
\setlist[itemize,1]{label={\color{accent}\textbullet}}
\usepackage{graphicx}
\usepackage{caption}
\captionsetup{font=small,labelfont={bf,sf,color=accent},textfont=sf,labelsep=quad,
  skip=6pt}
\captionsetup[table]{position=top}
\usepackage{float}
\usepackage{booktabs}
\usepackage{array}
\usepackage{tabularx}
\usepackage{longtable}
\usepackage{xltabular}
\usepackage{multirow}
\usepackage{etoolbox}
\newcolumntype{L}{>{\raggedright\arraybackslash}X}       % ragged-right X column
\newcolumntype{P}[1]{>{\raggedright\arraybackslash}p{#1}} % ragged-right fixed-width column
\renewcommand{\arraystretch}{1.2}
\setlength{\aboverulesep}{0.5ex}
\setlength{\belowrulesep}{0.6ex}
\setlength{\heavyrulewidth}{0.9pt}
\setlength{\lightrulewidth}{0.45pt}
\AtBeginEnvironment{tabularx}{\small}
\AtBeginEnvironment{xltabular}{\small}
\newcommand{\tabhead}{\leavevmode\sffamily\bfseries\color{titlecolor}}   % the header cells of a table

% ---- Code --------------------------------------------------------------------------------
\usepackage{listings}
\lstdefinelanguage{Rust}{
  morekeywords={as,break,const,continue,crate,else,enum,extern,false,fn,for,if,impl,in,let,
    loop,match,mod,move,mut,pub,ref,return,self,Self,static,struct,super,trait,true,type,
    unsafe,use,where,while,dyn,Some,None,Ok,Err,Option,Result,Vec,String,u8,u16,u32,u64,
    i16,i32,i64,f32,f64,usize,isize,bool,str},
  sensitive=true,
  morecomment=[l]{//},
  morecomment=[s]{/*}{*/},
  morestring=[b]",
}
\lstdefinelanguage{TOML}{
  morecomment=[l]{\#},
  morestring=[b]",
  morestring=[b]',
  morekeywords={true,false},
}
\lstdefinelanguage{CMake}{
  morekeywords={add_subdirectory,target_link_libraries,set,cmake_minimum_required,project,
    add_executable,find_package,PRIVATE,PUBLIC,CACHE,BOOL,FORCE,OFF,ON,EXCLUDE_FROM_ALL},
  sensitive=false,
  morecomment=[l]{\#},
  morestring=[b]",
}
\lstdefinelanguage{Julia}{
  morekeywords={abstract,baremodule,begin,break,catch,const,continue,do,else,elseif,end,
    export,false,finally,for,function,global,if,import,let,local,macro,module,mutable,
    nothing,quote,return,struct,true,try,using,while,where,missing},
  sensitive=true,
  morecomment=[l]\#,
  morecomment=[n]{\#=}{=\#},
  morestring=[b]",
}
\lstdefinelanguage{Kotlin}{
  morekeywords={as,break,class,continue,do,else,false,for,fun,if,in,interface,is,null,
    object,package,return,super,this,throw,true,try,typealias,val,var,when,while,by,
    catch,constructor,data,enum,finally,import,internal,private,public,vararg,
    Int,Long,Double,String,Boolean,List,Map,Pair},
  sensitive=true,
  morecomment=[l]{//},
  morecomment=[s]{/*}{*/},
  morestring=[b]",
}
\lstdefinelanguage{JSON}{
  morestring=[b]",
  morekeywords={true,false,null},
}
\lstset{
  basicstyle=\ttfamily\small,
  columns=fullflexible,
  keepspaces=true,
  upquote=true,
  showstringspaces=false,
  breaklines=true,
  breakatwhitespace=false,
  postbreak=\mbox{\textcolor{codecomment}{$\hookrightarrow$}\space},
  backgroundcolor=\color{codebg},
  frame=l,
  framerule=2pt,
  rulecolor=\color{okskyblue},
  framesep=6pt,
  xleftmargin=8pt,
  xrightmargin=0pt,
  aboveskip=0.8em,
  belowskip=0.6em,
  tabsize=4,
  keywordstyle=\color{codekeyword}\bfseries,
  commentstyle=\color{codecomment}\itshape,
  stringstyle=\color{codestring},
}
\lstdefinestyle{text}{language={},keywordstyle={},commentstyle={},stringstyle={}}
% The styles of the listings of the reference chapters. docs/user-manual/check_manual.py
% finds them by these names: "signature" is the call syntax of a function (one listing
% right after each \apientry), "example" is an example that the check runs (Python), and
% "shell" is a command line.
\lstdefinestyle{signature}{backgroundcolor=\color{white},rulecolor=\color{accent},
  basicstyle=\ttfamily\small\bfseries,keywordstyle=\color{codekeyword}}
\lstdefinestyle{example}{}
\lstdefinestyle{shell}{language=bash,keywordstyle={},commentstyle=\color{codecomment},
  stringstyle={}}

% ---- Boxes -------------------------------------------------------------------------------
\usepackage{tikz}
\usepackage[most]{tcolorbox}
\tcbset{wfdbxbox/.style={enhanced,breakable,boxrule=0pt,leftrule=3pt,arc=0pt,outer arc=0pt,
  left=8pt,right=8pt,top=5pt,bottom=5pt,before skip=0.8em,after skip=0.8em,
  fonttitle=\sffamily\bfseries\small,attach title to upper={\quad}}}
\newtcolorbox{note}[1][Note]{wfdbxbox,colback=noteback,colframe=accent,coltitle=accent,
  title={#1}}
\newtcolorbox{tip}[1][Tip]{wfdbxbox,colback=tipback,colframe=okgreen,coltitle=okgreen,
  title={#1}}
\newtcolorbox{phibox}[1][Patient data]{wfdbxbox,colback=phiback,colframe=okviolet,
  coltitle=okviolet,title={#1}}

% ---- Links and references ----------------------------------------------------------------
\usepackage{xurl}
\usepackage[bookmarksnumbered,bookmarksopen,bookmarksopenlevel=0]{hyperref}
\hypersetup{colorlinks=true,linkcolor=accent,urlcolor=accent,citecolor=accent,
  pdftitle={WFDB-x: User's manual},
  pdfauthor={Reza Sameni},
  pdfsubject={The WFDB Software Package in Python, C++, Rust, Julia, R, MATLAB/Octave
    and Kotlin},
  pdfkeywords={WFDB, PhysioNet, physiological signals, ECG, annotations, Python, C++,
    Rust, Julia, R, MATLAB, Octave, Kotlin}}
\usepackage[nameinlink,capitalise,noabbrev]{cleveref}

% ---- Manual macros -----------------------------------------------------------------------
\newcommand{\wfdbxversion}{1.2.1}      % the version of WFDB-x described
\newcommand{\wfdbxdate}{November 2026}
\newcommand{\wfdbversion}{10.7.0}     % the version of the WFDB Software Package implemented
\newcommand{\wfdbx}{WFDB-x}

% An underscore that allows a line break after it, so long identifiers wrap in tables.
\DeclareRobustCommand{\_}{\ifmmode\text{\textunderscore}\else\textunderscore\allowbreak\fi}

\newcommand{\code}[1]{\texttt{#1}}          % inline code, identifiers, values (escape _ # % & { })
\newcommand{\cmd}[1]{\texttt{#1}}           % a command line or a subcommand
\newcommand{\opt}[1]{\texttt{#1}}           % a command-line option
% A file or folder of the repository or an output: may break after each /.
\makeatletter
\DeclareRobustCommand{\file}[1]{\texttt{\wfdbx@file#1/\@nil}}
\def\wfdbx@file#1/#2\@nil{#1\if\relax\detokenize{#2}\relax\else/\allowbreak\wfdbx@file#2\@nil\fi}
\makeatother
\newcommand{\phimark}{\textsuperscript{\textcolor{okviolet}{\sffamily\bfseries PHI}}}
\pdfstringdefDisableCommands{\def\_{\textunderscore}\def\allowbreak{}\def\code#1{#1}%
  \def\file#1{#1}\def\cmd#1{#1}\def\opt#1{#1}\def\phimark{}}

% A function or command entry of a reference chapter: \apientry{name}{one-line summary}
\newcommand{\apientry}[2]{\par\addvspace{0.9em}\noindent
  {\color{rulecolor}\rule{\linewidth}{0.4pt}}\par\nopagebreak\vspace{0.15em}\noindent
  {\large\ttfamily\bfseries\color{titlecolor}#1}\hfill{\small\sffamily\color{mutedtext}#2}\par
  \nopagebreak\smallskip}
% Parameter lists of reference entries: \begin{params} \item[\code{name}] meaning \end{params}
\newlist{params}{description}{1}
\setlist[params]{font=\normalfont\ttfamily\bfseries,leftmargin=2.2em,style=nextline,
  itemsep=0.1em}

% Return value and errors of a reference entry: \returns{...} \raises{...}
\newcommand{\returns}[1]{\par\noindent{\sffamily\bfseries\small\color{titlecolor}Returns}
  \quad #1\par}
\newcommand{\raises}[1]{\par\noindent{\sffamily\bfseries\small\color{titlecolor}Errors}
  \quad #1\par}
% The language of a chapter, for the running heads and the check (\chapterlanguage{python})
\newcommand{\chapterlanguage}[1]{}

